
%%%%%%%%%%%%%%%%%%%%%%%%%%%%%%%%%%%%%%%%%%%%%%%%%%%%%%%%%%%
%% Inicio de la introducción
%%%%%%%%%%%%%%%%%%%%%%%%%%%%%%%%%%%%%%%%%%%%%%%%%%%%%%%%%%%

\chapter{Introducción}
%%---------------------------------------------------------

En el contexto actual de la transformación digital, el modelado de procesos de negocio se ha convertido en una herramienta esencial para las organizaciones que desean comprender, documentar y mejorar la forma en que entregan valor a sus clientes. El estándar \textit{Business Process Model and Notation} (BPMN) es ampliamente utilizado para representar estos procesos de forma gráfica y estructurada. Sin embargo, la creación manual de modelos BPMN requiere conocimientos especializados y supone un esfuerzo considerable, especialmente cuando se trata de procesos complejos o cuando se desea incorporar un análisis del valor percibido por el cliente.

En este contexto, los modelos de lenguaje de gran escala (\textit{Large Language Models}, LLMs) ofrecen nuevas posibilidades para automatizar la generación y el análisis de modelos de procesos. Su capacidad para interpretar descripciones en lenguaje natural y generar artefactos estructurados los convierte en aliados prometedores para asistir a analistas y profesionales en la construcción de modelos BPMN.

Este Trabajo de Fin de Máster (TFM) aborda el desarrollo de una extensión de la herramienta bpmn.io que integra un LLM para la generación y análisis automático de modelos BPMN a partir de descripciones textuales o conversacionales. La herramienta permite, de forma iterativa, construir modelos de procesos de negocio e identificar los valores asociados al cliente según la clasificación PERVAL. Para validar su utilidad, se diseña y ejecuta un estudio empírico en el que se comparan los resultados de la herramienta con los de un analista experto.

%%---------------------------------------------------------
\section{Motivación}
%%---------------------------------------------------------

La motivación de este trabajo surge de la observación de un doble desafío en el ámbito del análisis de procesos de negocio orientados al valor. Por un lado, la construcción de modelos BPMN de calidad es una tarea técnicamente exigente que requiere conocimiento del estándar, comprensión del dominio y experiencia en análisis de procesos. Por otro lado, los procesos de comercio electrónico están diseñados para entregar valor al cliente mediante atributos como la calidad, el precio y los elementos emocionales y sociales, pero no siempre el valor modelado por la organización corresponde al valor percibido por los usuarios durante la interacción real con el proceso.

Este desalineamiento entre el valor diseñado y el valor percibido puede tener un impacto directo en la experiencia del cliente y en la eficacia del proceso. Sin embargo, su identificación y corrección requieren de un análisis sistemático que, en la práctica, suele ser manual y costoso en tiempo y recursos.

Los avances recientes en LLMs ofrecen una oportunidad para abordar estos desafíos de forma más eficiente. Su capacidad para comprender lenguaje natural y generar artefactos estructurados los convierte en candidatos idóneos para asistir en la generación de modelos BPMN y en la identificación del valor para el cliente.

La realización de este TFM ha permitido aplicar de forma práctica los conocimientos adquiridos durante el máster universitario, en especial en las áreas de ingeniería del software, modelado de procesos y tecnologías de inteligencia artificial.

%%---------------------------------------------------------
\section{Objetivos}
%%---------------------------------------------------------

El objetivo principal de este trabajo es desarrollar una herramienta basada en LLMs que permita generar y analizar modelos BPMN de forma automática a partir de descripciones textuales o conversacionales, integrada como extensión sobre una versión local de bpmn.io, y validar su utilidad mediante un estudio empírico.

A partir de este objetivo principal, se definen los siguientes objetivos secundarios:

\begin{itemize}
    \item[$\bullet$] diseñar e implementar una extensión funcional sobre bpmn.io que permita la generación iterativa de modelos BPMN mediante interacción en lenguaje natural con un LLM,
    \item[$\bullet$] incorporar en la herramienta la capacidad de identificar y clasificar los valores asociados al cliente según la clasificación PERVAL,
    \item[$\bullet$] diseñar y ejecutar un estudio empírico que permita comparar los modelos BPMN y los valores identificados por la herramienta con los generados por un analista experto de forma manual,
    \item[$\bullet$] analizar y evaluar el grado de alineamiento entre los resultados de la herramienta y el criterio humano experto tanto en la construcción del modelo como en la identificación del valor, y
    \item[$\bullet$] extraer conclusiones sobre la viabilidad y las limitaciones del uso de LLMs para el modelado de procesos de negocio orientado al valor.
\end{itemize}

%%---------------------------------------------------------
\section{Impacto Esperado}
%%---------------------------------------------------------

El impacto esperado de este trabajo es doble. En primer lugar, desde un punto de vista tecnológico, se espera demostrar la viabilidad de integrar LLMs en herramientas de modelado de procesos de negocio estándar, abriendo una línea de trabajo orientada a la asistencia inteligente en el modelado BPMN. En segundo lugar, desde un punto de vista metodológico, el estudio empírico aportará evidencia sobre el grado en que las herramientas basadas en LLMs pueden aproximarse al juicio de un analista humano experto, tanto en la calidad del modelo generado como en la identificación del valor para el cliente.

Este proyecto se alinea con los siguientes Objetivos de Desarrollo Sostenible (ODS):

\begin{itemize}
    \item[$\bullet$] \textbf{ODS 9: industria, innovación e infraestructura.} El uso de LLMs para la automatización del modelado de procesos contribuye al avance de soluciones tecnológicas innovadoras aplicadas a la ingeniería del software y a la mejora de las herramientas de análisis empresarial.
    \item[$\bullet$] \textbf{ODS 12: producción y consumo responsables.} Al facilitar la identificación del valor percibido por el cliente en los procesos de comercio electrónico, la herramienta desarrollada contribuye a un diseño de procesos más eficiente y ajustado a las necesidades reales del usuario, reduciendo el desperdicio de recursos en actividades que no generan valor.
\end{itemize}

%%---------------------------------------------------------
\section{Metodología}
%%---------------------------------------------------------

La metodología empleada en este trabajo combina el desarrollo de una herramienta software con el diseño y ejecución de un estudio empírico. Para la construcción de la herramienta se parte de una revisión de la literatura en torno a los conceptos y tecnologías clave implicados. Su utilidad se valida posteriormente mediante un estudio empírico comparativo, en el que los modelos BPMN y los valores identificados automáticamente por la herramienta se contrastan con los generados por un analista experto humano a partir del mismo caso de estudio: un proceso de compra en línea inspirado en escenarios reales de comercio electrónico. Este enfoque permite evaluar hasta qué punto la herramienta es capaz de aproximarse al criterio humano tanto en la construcción del modelo como en la identificación del valor.

%%---------------------------------------------------------
\section{Estructura}
%%---------------------------------------------------------

La estructura de este TFM ha sido diseñada para guiar al lector a lo largo del proceso de investigación y desarrollo llevado a cabo. A continuación, se describe el contenido de cada capítulo y los anexos.

\begin{itemize}
    \item[$\bullet$] \textbf{Introducción:} en este capítulo se contextualiza el proyecto, explicando la motivación que lo origina, los objetivos planteados y el impacto esperado del trabajo. También se describe la metodología de investigación seguida, la estructura del documento y las convenciones tipográficas adoptadas.

    \item[$\bullet$] \textbf{Marco teórico y estado del arte:} en este capítulo se revisan los fundamentos teóricos y el contexto tecnológico necesarios para comprender el trabajo. Se estudia el estándar BPMN y su aplicación al modelado de procesos de negocio orientados al valor, se analizan las capacidades de los LLMs para la generación y el análisis de artefactos software a partir de lenguaje natural, y se presenta la clasificación PERVAL como marco para identificar y categorizar el valor percibido por el cliente. Asimismo, se describe la herramienta bpmn.io como plataforma de modelado sobre la que se desarrolla la extensión, y se revisan trabajos relacionados con la generación automática de modelos de procesos y el uso de inteligencia artificial en ingeniería del software.

    \item[$\bullet$] \textbf{Diseño e implementación de la herramienta:} en este capítulo se presenta la solución desarrollada. Se describe la arquitectura general de la extensión sobre bpmn.io, detallando los componentes que la integran y la forma en que interactúan entre sí. Se explica el módulo de generación iterativa de modelos BPMN mediante un LLM a partir de descripciones textuales o conversacionales, así como el módulo de identificación y clasificación automática del valor del cliente según la escala PERVAL. También se justifican las decisiones de diseño adoptadas y se describen las tecnologías y herramientas empleadas en el desarrollo.

    \item[$\bullet$] \textbf{Estudio empírico:} en este capítulo se presenta el diseño y la ejecución del estudio empírico llevado a cabo para evaluar la herramienta desarrollada. Se define el caso de estudio utilizado, basado en un proceso de compra en línea inspirado en escenarios reales de comercio electrónico, y se describe el procedimiento seguido para las iteraciones de modelado y análisis de valor. Se detallan los criterios y métricas de comparación empleados para contrastar los modelos BPMN y los valores identificados automáticamente con los generados por un analista experto humano de forma manual.

    \item[$\bullet$] \textbf{Resultados:} en este capítulo se presentan y analizan los resultados obtenidos en el estudio empírico. Se comparan los modelos BPMN generados por la herramienta con los elaborados por el analista experto, así como los valores identificados automáticamente con los clasificados manualmente según PERVAL. Se discute el grado de alineamiento entre ambos enfoques, se interpretan las diferencias encontradas y se evalúa hasta qué punto la herramienta es capaz de aproximarse al criterio humano.

    \item[$\bullet$] \textbf{Conclusiones:} en este capítulo se recogen las principales conclusiones del trabajo, valorando en qué medida se han alcanzado los objetivos planteados. Se reflexiona sobre las limitaciones de la herramienta y del estudio empírico, se analiza la relación del trabajo desarrollado con la formación académica recibida y se proponen posibles líneas de trabajo futuro que permitan ampliar o mejorar los resultados obtenidos.

    \item[$\bullet$] \textbf{Bibliografía:} se recopilan todas las fuentes de información consultadas durante la realización del TFM, incluyendo libros, artículos de investigación, publicaciones académicas, documentación técnica y recursos digitales que han sido relevantes para el desarrollo del proyecto.

    \item[$\bullet$] \textbf{Anexos:} contienen información complementaria al cuerpo principal del documento, incluyendo material de apoyo al estudio empírico, datos adicionales de la evaluación y otros recursos que facilitan la comprensión y la reproducibilidad del trabajo realizado.
\end{itemize}

%%---------------------------------------------------------
\section{Convenciones}
%%---------------------------------------------------------

En este TFM se adoptará una convención específica para el marcado de términos en lenguas extranjeras. La primera vez que aparezca un término técnico o una palabra que no pertenezca al idioma principal del documento, se destacará en cursiva. Posteriormente, podrá utilizarse sin cursiva si ya ha sido introducido previamente.

%%%%%%%%%%%%%%%%%%%%%%%%%%%%%%%%%%%%%%%%%%%%%%%%%%%%%%%%%%%
%% Final de la introducción
%%%%%%%%%%%%%%%%%%%%%%%%%%%%%%%%%%%%%%%%%%%%%%%%%%%%%%%%%%%
